%================================================================
% 作業ログ テンプレート
%   教員配布の「作業ログ_配布.pdf」と同じ章構成・要素・階層を LaTeX で再現。
%   実際の作業ログを書くときは本ファイルをコピーして書き換える。
%
%   ビルド:
%     cd report/作業ログテンプレート
%     docker run --rm -v "$(pwd):/workspace" -w /workspace \
%       ghcr.io/paperist/texlive-ja:debian latexmk 作業ログテンプレート.tex
%   出力: 作業ログテンプレート.pdf（配布原本とは別ファイル）
%================================================================
\documentclass[uplatex,dvipdfmx,a4paper,11pt]{jsarticle}

% --- 余白 ---
\usepackage[top=25mm,bottom=25mm,left=25mm,right=25mm]{geometry}

% --- 表・箇条書き・図 ---
\usepackage{array}
\usepackage{tabularx}
\usepackage{booktabs}
\usepackage{enumitem}
\usepackage{graphicx}

% --- 段落間隔（Wordのような見た目に寄せる）---
\setlength{\parindent}{0pt}
\setlength{\parskip}{0.6em}

% --- セクション見出しを「番号. 太字」+ 章末に水平線を入れる ---
\usepackage{titlesec}
\titleformat{\section}
  {\normalfont\Large\bfseries}{\thesection.}{0.6em}{}
\titleformat{\subsection}
  {\normalfont\large\bfseries}{\thesubsection}{0.6em}{}
\titlespacing*{\section}{0pt}{1.2em}{0.6em}
\titlespacing*{\subsection}{0pt}{1.0em}{0.4em}

% --- 章末水平線（手で \sectionrule と書いて区切る）---
\newcommand{\sectionrule}{\par\medskip\noindent\rule{\linewidth}{0.6pt}\par\medskip}

% --- 箇条書きの記号を Word 風（・→○→数字）に揃える ---
\renewcommand{\labelitemi}{\textbullet}
\renewcommand{\labelitemii}{\textopenbullet}
\setlist[itemize]{leftmargin=1.6em,itemsep=0.1em,topsep=0.2em}
\setlist[enumerate]{leftmargin=2.0em,itemsep=0.1em,topsep=0.2em}

% \textopenbullet が無い環境向けのフォールバック（jsarticle標準では○が出る）
\providecommand{\textopenbullet}{\ensuremath{\circ}}

\begin{document}

%================================================================
% ヘッダ
%================================================================
{\LARGE\bfseries 作業ログ}\par\bigskip

報告者：[氏名]

報告日：[日付]

プロジェクト名：[プロジェクト名]

\sectionrule

%================================================================
% 1. 進捗サマリー
%================================================================
\section{進捗サマリー（要約）}

\begin{itemize}
  \item 全体進捗率：[％]
  \item マイルストーン達成状況：[達成／遅延／前倒し]
  \item 主要成果：
    \begin{itemize}
      \item {}[成果1]
      \item {}[成果2]
    \end{itemize}
  \item 重大課題（影響度：高）：
    \begin{itemize}
      \item {}[課題概要 + 一言で影響]
    \end{itemize}
\end{itemize}

\sectionrule

%================================================================
% 2. 作業ログ（詳細トラッキング）
%================================================================
\section{作業ログ（詳細トラッキング）}

\begin{tabularx}{\linewidth}{@{}lllXlXX@{}}
\toprule
作業項目 & 開始日 & 予定完了日 & 状態 & 進捗率 & 依存関係・ブロッカー & コメント \\
\midrule
作業A & [日付] & [日付] & 未着手／進行中／レビュー中／完了 & [％] & [有無・内容] & [特記事項] \\
作業B & [日付] & [日付] &                                  &      &              &            \\
\bottomrule
\end{tabularx}

\sectionrule

%================================================================
% 3. 課題管理
%================================================================
\section{課題管理}

\begin{tabularx}{\linewidth}{@{}lXllXll@{}}
\toprule
課題 & 発生原因 & 影響度 & 優先度 & 対策 & 期限 & 状態 \\
\midrule
課題1 & [原因] & 高／中／低 & 高／中／低 & [対策案] & [日付] & 未対応／対応中／解決 \\
\bottomrule
\end{tabularx}

\medskip
※影響度＝プロジェクトへのインパクト

※優先度＝対応の緊急性

\sectionrule

%================================================================
% 4. AI利用状況と検証
%================================================================
\section{AI利用状況と検証（再現性重視）}

\subsection{利用内容}

\begin{itemize}
  \item 利用目的：[探索／設計／実装／レビュー 等]
  \item 入力（プロンプト要旨）：
    \begin{itemize}
      \item {}[何をどのように聞いたか]
    \end{itemize}
  \item 出力概要：
    \begin{itemize}
      \item {}[何が得られたか]
    \end{itemize}
\end{itemize}

\sectionrule

\subsection{検証方法}

※第三者が再現できるレベルで記述

\begin{itemize}
  \item 検証手法：
    \begin{itemize}
      \item {}[実験／比較／理論確認など]
    \end{itemize}
  \item 参照資料：
    \begin{itemize}
      \item {}[論文・書籍・公式ドキュメント等]
    \end{itemize}
  \item 検証手順：
    \begin{enumerate}
      \item {}[手順1]
      \item {}[手順2]
    \end{enumerate}
\end{itemize}

\sectionrule

\subsection{ファクトチェック結果}

\begin{itemize}
  \item 一致点：
    \begin{itemize}
      \item {}[既存知識と整合している部分]
    \end{itemize}
  \item 不一致・疑義：
    \begin{itemize}
      \item {}[差異の内容]
    \end{itemize}
  \item 最終判断：
    \begin{itemize}
      \item 採用／保留／棄却
    \end{itemize}
  \item 根拠：
    \begin{itemize}
      \item {}[資料・理由]
    \end{itemize}
\end{itemize}

\sectionrule

%================================================================
% 5. 次回計画
%================================================================
\section{次回計画}

\begin{itemize}
  \item 目標：
    \begin{itemize}
      \item {}[タスク内容]
    \end{itemize}
  \item 達成基準：
    \begin{itemize}
      \item {}[検証可能な条件]
    \end{itemize}
  \item 期限：
    \begin{itemize}
      \item {}[日付]
    \end{itemize}
\end{itemize}

\sectionrule

%================================================================
% 6. その他
%================================================================
\section{その他}

\begin{itemize}
  \item {}[連絡事項・懸念点など]
\end{itemize}

\sectionrule

%================================================================
% 7. 付録
%================================================================
\section{付録}

\begin{itemize}
  \item 添付資料：
    \begin{itemize}
      \item ソースコード：[ファイル or リポジトリURL]
      \item 回路図／設計図：[ファイル]
    \end{itemize}
  \item 添付範囲：
    \begin{itemize}
      \item 全体／差分（どちらか明記）
    \end{itemize}
\end{itemize}

\sectionrule

\end{document}
